\documentclass[12pt,a4paper]{article}
\usepackage[utf8]{inputenc}
\usepackage[T5]{fontenc}
\usepackage[vietnamese]{babel}
\usepackage{graphicx}
\usepackage{tabularx}
\usepackage{booktabs}
\usepackage{geometry}
\usepackage{hyperref}

% Căn lề chuẩn đồ án
\geometry{a4paper, left=3cm, right=2cm, top=2.5cm, bottom=2.5cm}

\title{\textbf{CHƯƠNG 3: THIẾT KẾ CƠ SỞ DỮ LIỆU} \\ \large{HỆ THỐNG KẾ HOẠCH TÀI CHÍNH (KHTC)}}
\author{Hoàng Mạnh Kiên}
\date{Tháng 3 Năm 2026}

\begin{document}

\maketitle

\section{Giới thiệu chung về Cơ sở dữ liệu}
Hệ thống Kế hoạch Tài chính (KHTC) được thiết kế theo kiến trúc điều khiển bằng siêu dữ liệu (Metadata-Driven). Hệ thống sử dụng mô hình cơ sở dữ liệu lai (Hybrid Database) kết hợp giữa SQL Server (để lưu trữ cấu trúc, phân quyền, quy trình), MongoDB (lưu trữ JSON linh hoạt từ các Form động) và ClickHouse (kho dữ liệu OLAP phục vụ truy vấn báo cáo tốc độ cao).

Cơ sở dữ liệu được chia thành 5 phân hệ cốt lõi để đảm bảo tính module hóa và dễ dàng mở rộng.

\section{Thiết kế chi tiết các bảng dữ liệu}

% ==============================
% PHÂN HỆ 1
% ==============================
\subsection{Phân hệ 1: Quản trị Hệ thống và Phân quyền (RBAC)}
Phân hệ này quản lý định danh người dùng, cấu trúc cây chức năng và ma trận phân quyền chi tiết.

\begin{table}[h!]
    \centering
    \caption{Bảng SYS\_USER (Quản lý người dùng)}
    \begin{tabularx}{\textwidth}{|l|l|l|X|}
        \hline
        \textbf{Trường dữ liệu} & \textbf{Kiểu dữ liệu} & \textbf{Ràng buộc} & \textbf{Mô tả chi tiết} \\ \hline
        UserID & INT & PK & Khóa chính, ID định danh người dùng. \\ \hline
        Username & VARCHAR(50) & Unique & Tên đăng nhập (VD: cuong.vp). \\ \hline
        FullName & NVARCHAR(200) & & Họ và tên đầy đủ của người dùng. \\ \hline
        EntityCode & VARCHAR(50) & FK & Mã đơn vị trực thuộc (Link với DIM\_ENTITY). \\ \hline
        IsActive & BIT & & Trạng thái hoạt động (1: Active, 0: Khóa). \\ \hline
    \end{tabularx}
\end{table}

\begin{table}[h!]
    \centering
    \caption{Bảng SYS\_ROLE và SYS\_USER\_ROLE (Phân quyền vai trò)}
    \begin{tabularx}{\textwidth}{|l|l|l|X|}
        \hline
        \textbf{Bảng} & \textbf{Trường dữ liệu} & \textbf{Kiểu dữ liệu} & \textbf{Mô tả chi tiết} \\ \hline
        SYS\_ROLE & RoleID & INT (PK) & Khóa chính, ID vai trò. \\ \cline{2-4}
                  & RoleName & NVARCHAR(100) & Tên vai trò (Admin, Ban KH, Kế toán). \\ \hline
        SYS\_USER\_ROLE & UserID & INT (FK) & Liên kết tới người dùng. \\ \cline{2-4}
                  & RoleID & INT (FK) & Liên kết tới vai trò tương ứng. \\ \hline
    \end{tabularx}
\end{table}

\begin{table}[h!]
    \centering
    \caption{Bảng SYS\_MENU và SYS\_ROLE\_PERMISSION (Cây tính năng)}
    \begin{tabularx}{\textwidth}{|l|l|X|}
        \hline
        \textbf{Bảng} & \textbf{Trường dữ liệu} & \textbf{Mô tả chi tiết} \\ \hline
        SYS\_MENU & MenuID (PK) & Khóa chính, ID Menu. \\ \cline{2-3}
                  & MenuName & Tên hiển thị Menu (VD: Kế hoạch sản xuất). \\ \cline{2-3}
                  & ParentID, Url & Menu cha và đường dẫn truy cập App. \\ \cline{2-3}
                  & FormID & Mã Form liên kết (VD: F01\_SXKD). \\ \hline
        SYS\_ROLE\_ & RoleID (FK) & Liên kết với Role. \\ \cline{2-3}
        PERMISSION& MenuID (FK) & Liên kết với Menu. \\ \cline{2-3}
                  & CanRead, CanWrite & Quyền xem và quyền sửa/nhập liệu (BIT). \\ \hline
    \end{tabularx}
\end{table}

% ==============================
% PHÂN HỆ 2
% ==============================
\subsection{Phân hệ 2: Danh mục Cốt lõi (Master Data)}
Cung cấp bộ từ điển chuẩn hóa để tổng hợp số liệu từ các đơn vị thành viên lên Tập đoàn.

\begin{table}[h!]
    \centering
    \caption{Bảng DIM\_ACCOUNT (Danh mục Chỉ tiêu - Cây báo cáo)}
    \begin{tabularx}{\textwidth}{|l|l|l|X|}
        \hline
        \textbf{Trường dữ liệu} & \textbf{Kiểu dữ liệu} & \textbf{Ràng buộc} & \textbf{Mô tả chi tiết} \\ \hline
        AccountID & INT & PK, Identity & Khóa chính tự tăng. \\ \hline
        AccountCode & VARCHAR(50) & Unique & Mã chỉ tiêu chuẩn (VD: DOANH\_THU). \\ \hline
        AccountName & NVARCHAR(255) & & Tên chỉ tiêu. \\ \hline
        ParentID & INT & FK & Chỉ tiêu cha. \\ \hline
        AccountType & VARCHAR(20) & & Phân loại: REVENUE, EXPENSE, ASSET. \\ \hline
        DataStorage & VARCHAR(20) & & Cách lưu: STORE, DYNAMIC\_CALC, LABEL. \\ \hline
        Operator & INT & Default 1 & Dấu tính toán khi cộng lên cha (+1, -1). \\ \hline
        Formula & NVARCHAR(MAX) & & Công thức dọc (VD: [LN\_TruocThue] - [Thue]). \\ \hline
        Unit & NVARCHAR(50) & & Đơn vị tính (Tỷ đồng, kWh). \\ \hline
        IsLeaf & BIT & & 1: Chỉ tiêu lá, 0: Chỉ tiêu tổng hợp. \\ \hline
    \end{tabularx}
\end{table}

\begin{table}[h!]
    \centering
    \caption{Bảng DIM\_ENTITY và DIM\_VERSION (Danh mục Đơn vị và Phiên bản)}
    \begin{tabularx}{\textwidth}{|l|l|l|X|}
        \hline
        \textbf{Bảng} & \textbf{Trường dữ liệu} & \textbf{Ràng buộc} & \textbf{Mô tả chi tiết} \\ \hline
        DIM\_ENTITY & EntityID, EntityCode & PK, Unique & ID và Mã đơn vị (VD: EVN, GENCO1). \\ \cline{2-4}
                    & EntityName, ParentID & & Tên đơn vị và ID đơn vị cha. \\ \cline{2-4}
                    & EntityType, Level & & Phân loại đơn vị (HOLDING, PC) và cấp độ. \\ \cline{2-4}
                    & Path & & Cây phả hệ (VD: /1/5/12/) để truy vấn nhanh. \\ \hline
        DIM\_VERSION& VersionCode & PK & Mã phiên bản số liệu (V1, Final). \\ \cline{2-4}
                    & VersionType, IsLocked& & Nhóm phiên bản và Trạng thái khóa số liệu. \\ \hline
    \end{tabularx}
\end{table}

% ==============================
% PHÂN HỆ 3
% ==============================
\subsection{Phân hệ 3: Cấu hình Mẫu biểu và Ánh xạ (Dynamic Forms)}
Nhóm bảng định nghĩa cấu trúc giao diện Handsontable và thuật toán ánh xạ tọa độ lưới với cơ sở dữ liệu.

\begin{table}[h!]
    \centering
    \caption{Bảng SYS\_FORM\_TEMPLATE và SYS\_FORM\_VERSION (Khung giao diện)}
    \begin{tabularx}{\textwidth}{|l|l|X|}
        \hline
        \textbf{Bảng} & \textbf{Trường dữ liệu} & \textbf{Mô tả chi tiết} \\ \hline
        SYS\_FORM\_ & FormID (PK) & Mã loại báo cáo (VD: F01\_SXKD). \\ \cline{2-3}
        TEMPLATE  & FormName, OrgList & Tên báo cáo và Loại đơn vị áp dụng. \\ \cline{2-3}
                  & LayoutConfig & Cấu trúc JSON chung của Form. \\ \cline{2-3}
                  & IsDynamicRow & 1: Cho phép thêm dòng, 0: Form tĩnh. \\ \hline
        SYS\_FORM\_ & VersionID (PK) & ID bản mẫu theo thời gian. \\ \cline{2-3}
        VERSION   & FormID, Year & Liên kết Template gốc và Năm áp dụng. \\ \cline{2-3}
                  & \textbf{LayoutJSON} & \textbf{Cấu hình Handsontable hoàn chỉnh (Cột, gộp ô).} \\ \hline
    \end{tabularx}
\end{table}

\begin{table}[h!]
    \centering
    \caption{Bảng SYS\_FORM\_MAPPING (Ánh xạ tọa độ Metadata)}
    \begin{tabularx}{\textwidth}{|l|l|l|X|}
        \hline
        \textbf{Trường dữ liệu} & \textbf{Kiểu dữ liệu} & \textbf{Ràng buộc} & \textbf{Mô tả chi tiết} \\ \hline
        MappingID & BIGINT & PK & Khóa chính. \\ \hline
        VersionID & INT & FK & Thuộc bản mẫu form nào. \\ \hline
        RowKey & VARCHAR(50) & Index & UUID của dòng, sợi dây neo giữ chỉ tiêu. \\ \hline
        ColKey & VARCHAR(50) & Index & Định danh cột (VD: C\_JAN, C\_Q1). \\ \hline
        AccountCode & VARCHAR(50) & FK & \textbf{Mã chỉ tiêu tài chính liên kết tại tọa độ này.} \\ \hline
        IsReadOnly & BIT & & Ghi đè cấu hình, ép buộc ô này chỉ đọc. \\ \hline
        Formula & NVARCHAR(MAX)& & Công thức HyperFormula ngang riêng cho ô. \\ \hline
        Style\_JSON & NVARCHAR(MAX)& & Định dạng riêng (màu sắc, border) cho ô. \\ \hline
    \end{tabularx}
\end{table}

% ==============================
% PHÂN HỆ 4
% ==============================
\subsection{Phân hệ 4: Lưu trữ Dữ liệu Giao dịch và Phân tích}
Nơi chứa các con số nghiệp vụ thực tế (Fact Data).

\begin{table}[h!]
    \centering
    \caption{Bảng FACT\_PLANNING\_DATA (Lưu trữ Transactional SQL)}
    \begin{tabularx}{\textwidth}{|l|l|l|X|}
        \hline
        \textbf{Trường dữ liệu} & \textbf{Kiểu dữ liệu} & \textbf{Ràng buộc} & \textbf{Mô tả chi tiết} \\ \hline
        OrgCode & VARCHAR(50) & PK, FK & Mã đơn vị sở hữu số liệu. \\ \hline
        AccountCode & VARCHAR(50) & PK, FK & Mã chỉ tiêu tài chính. \\ \hline
        Period & INT & PK & Kỳ dữ liệu (Tháng/Quý/Năm - VD: 202601). \\ \hline
        ScenarioCode & VARCHAR(20) & PK & Kịch bản (Plan, Actual, Forecast). \\ \hline
        VersionCode & VARCHAR(20) & PK, FK & Phiên bản số liệu. \\ \hline
        \textbf{Value} & DECIMAL(18,4)& & \textbf{Con số thực tế cuối cùng.} \\ \hline
        Currency, Year& VARCHAR, INT & & Tiền tệ và Năm dữ liệu (tối ưu Filter). \\ \hline
    \end{tabularx}
\end{table}

\begin{table}[h!]
    \centering
    \caption{Bảng Fact\_Financial\_Performance (Kho dữ liệu OLAP - ClickHouse)}
    \begin{tabularx}{\textwidth}{|l|l|X|}
        \hline
        \textbf{Trường dữ liệu} & \textbf{Kiểu ClickHouse} & \textbf{Mô tả chi tiết} \\ \hline
        EntityCode, AccountCode & String & Mã đơn vị và Mã chỉ tiêu đã làm phẳng. \\ \hline
        PeriodID, Version & UInt32, UInt16 & Kỳ dữ liệu và Phiên bản. \\ \hline
        Scenario & LowCardinality & Kịch bản dữ liệu. \\ \hline
        \textbf{Amount} & Decimal(18,4) & \textbf{Giá trị số liệu để tính tổng.} \\ \hline
        EntityPath & Array(String) & Mảng phả hệ đơn vị (VD: ['EVN', 'GENCO1']). \\ \hline
        AccountPath & Array(String) & Mảng phả hệ chỉ tiêu. Giúp Query không cần Join. \\ \hline
    \end{tabularx}
\end{table}

% ==============================
% PHÂN HỆ 5
% ==============================
\subsection{Phân hệ 5: Quy trình Phê duyệt và Ký số (Workflow)}

\begin{table}[h!]
    \centering
    \caption{Các bảng Cấu hình và Trạng thái Quy trình (Workflow)}
    \begin{tabularx}{\textwidth}{|l|l|X|}
        \hline
        \textbf{Bảng} & \textbf{Trường dữ liệu} & \textbf{Mô tả chi tiết} \\ \hline
        SYS\_WORKFLOW\_ & StepID, StepName & Thứ tự bước (1-10) và Tên bước (VD: Thẩm định). \\ \cline{2-3}
        CONFIG        & DeptCode, CanEdit& Phòng ban xử lý và Quyền sửa số. \\ \cline{2-3}
                      & RequiredPrevStep & Ràng buộc bước phải hoàn thành trước đó. \\ \hline
        FACT\_WORKFLOW\_& WfID (PK), EntityID& ID định danh tiến độ và Đơn vị thực hiện. \\ \cline{2-3}
        STATUS        & CurrentStepID & Đang đứng ở bước nào trong quy trình. \\ \cline{2-3}
                      & Status & Trạng thái (DRAFT, SUBMITTED, APPROVED). \\ \hline
        FACT\_        & SubmissionID (PK)& ID hồ sơ nộp. \\ \cline{2-3}
        SUBMISSION    & JsonData & Toàn bộ dữ liệu thô JSON lưu lại làm bằng chứng. \\ \hline
    \end{tabularx}
\end{table}

\begin{table}[h!]
    \centering
    \caption{Nhật ký Phê duyệt và Lưu trữ Ký số}
    \begin{tabularx}{\textwidth}{|l|l|X|}
        \hline
        \textbf{Bảng} & \textbf{Trường dữ liệu} & \textbf{Mô tả chi tiết} \\ \hline
        FACT\_APPROVAL\_& LogID (PK) & ID nhật ký thao tác. \\ \cline{2-3}
        HISTORY       & SubmissionID & Liên kết tới hồ sơ nộp. \\ \cline{2-3}
                      & Action, ActionBy & Hành động (APPROVE/REJECT) và Người thực hiện. \\ \cline{2-3}
                      & Comment, ActionTime& Ghi chú phê duyệt và Thời gian thực hiện. \\ \hline
        FACT\_DIGITAL\_ & SigID (PK) & ID chữ ký số. \\ \cline{2-3}
        SIGNATURE     & SubmissionID & Liên kết hồ sơ. \\ \cline{2-3}
                      & FilePath, FileHash & File PDF đã đóng mộc và Mã băm chống tráo đổi. \\ \cline{2-3}
                      & SignatureValue & Chuỗi chữ ký số mã hóa trả về từ Provider. \\ \hline
    \end{tabularx}
\end{table}

\section{Kết luận}
Kiến trúc cơ sở dữ liệu kết hợp đa nền tảng này đảm bảo hệ thống KHTC vừa có sự chặt chẽ về mặt quy trình nghiệp vụ (SQL Server), vừa có sự linh hoạt trong thiết kế biểu mẫu động (MongoDB), đồng thời tối ưu hóa được tốc độ kết xuất các báo cáo tài chính hợp nhất khổng lồ (ClickHouse).

\end{document}